\chapter{Related Work}
\label{sec:related-work}

\section{Distillation}

Distillation transfers behavior from a teacher model to a student model by training the student to match the teacher's predictions or generated outputs. The classical formulation was introduced as a way to compress the knowledge of a large model or ensemble into a smaller model using the teacher's softened output distribution \cite{hintonDistillingKnowledgeNeural2015}. In language modeling, the same principle has been adapted to large pretrained models \cite{tan2023GKDGeneralKnowledge}, large-scale pretraining systems \cite{wang2021ERNIE30Titan}, selective token-weighted objectives \cite{huang2025SelecTKDSelectiveTokenWeighted}, and cross-tokenizer teacher--student pairs \cite{boizardCrossTokenizerDistillationUniversal2025,minixhoferUniversalCrossTokenizerDistillation2025}. For reasoning models, sequence-level and chain-of-thought distillation are also central, with recent work showing that standard supervised fine-tuning pipelines can poorly represent the teacher's sequence-level distribution and proposing distribution-aligned sequence distillation for long-CoT reasoning \cite{yanDistributionAlignedSequenceDistillation2026}.

Most classical language-model distillation methods are \emph{off-policy}: the student is trained on a fixed data distribution rather than on rollouts generated by the current student. This fixed data can be human-written, sampled from the teacher, or generated by another system. The supervision can be hard-label supervised fine-tuning (SFT), where the student maximizes the likelihood of target tokens \cite{zotero-item-425}, or soft-label distillation, where the student matches teacher log-probabilities or full next-token distributions on the same fixed prefixes. In both cases, the training states are chosen before the student update rather than by the student's current policy.

Off-policy distillation is attractive because it is simple, stable, and easy to scale. Teacher completions and teacher log-probabilities can be generated once, cached, filtered, and reused across many training runs. This decouples expensive teacher inference from student optimization and makes off-policy distillation a practical baseline for large teachers. Its weakness is that the student is not trained on the prefixes it will actually visit at inference time. For autoregressive models, early student mistakes can move the model into states that are unlikely under the fixed dataset, creating exposure bias \cite{pmlr-v9-ross10a}.

On-policy distillation (OPD) addresses this mismatch by sampling rollouts from the student and asking the teacher to provide feedback at the student's own prefixes. The teacher no longer only produces complete demonstrations; it gives local supervision, often through next-token log-probabilities, on states induced by the student. Generalized Knowledge Distillation trains language models on self-generated sequences with teacher feedback and improves summarization, translation, arithmetic reasoning, and instruction tuning \cite{agarwalOnPolicyDistillationLanguage2024}. MiniLLM argues that reverse KL is better suited to on-policy generative distillation because it discourages the student from assigning high probability to tokens that the teacher considers unlikely \cite{guMiniLLMOnPolicyDistillation2026}. Recent OPD work further shows that useful training signals can be obtained even from sampled-token log-probabilities rather than full vocabulary distributions \cite{luOnPolicyDistillation2025}.

This thesis uses OPD as the main training technique because it lets the teacher supervise the student on the student's own reasoning trajectories. We nevertheless compare OPD against off-policy distillation methods, because off-policy SFT and teacher-generated traces remain the simplest and most common alternatives. This comparison is especially important in mathematical reasoning, where student errors often create prefixes that are absent from teacher-written solutions. More broadly, on-policy data can mitigate forgetting by preserving the model's own behavioral distribution during learning \cite{chenRetainingDoingRole2025}, and diagnostic work shows that on-policy distillation is most useful when the student rollout is wrong, while the best distillation context depends on both task and student capacity \cite{armandpourUnmaskingOnPolicyDistillation2026}.

\section{Addressing the capacity gap}
Distillation is usually motivated by the asymmetry between a larger, more capable teacher and a smaller, more efficient student. This asymmetry is also the source of a central difficulty: a stronger teacher is not always a better teacher for a much smaller student. Early studies \cite{mirzadehImprovedKnowledgeDistillation2019} show that student performance can degrade when the teacher is too large relative to the student, motivating intermediate teacher-assistant models to bridge the gap \cite{mirzadehImprovedKnowledgeDistillation2019}. Later work describes this failure mode as the ``curse of capacity gap,'' where a large teacher's additional capability does not transfer cleanly to a much smaller student \cite{zhangLiftingCurseCapacity2023}. This problem is central to our setting because the most appealing deployment regime is precisely the hardest one: using the strongest available teacher to train the smallest student that is still cheap and fast enough to serve.

Existing approaches to the capacity gap often introduce additional machinery to make the teacher signal easier for the student to learn. Teacher-assistant methods insert intermediate models between teacher and student \cite{mirzadehImprovedKnowledgeDistillation2019}. Progressive distillation guides the student along an optimized path from its initial distribution toward the teacher \cite{shi2021FollowYourPath}. MiniMoE changes the student architecture by adding minimal expert modules to improve transfer without greatly increasing inference cost \cite{zhangLiftingCurseCapacity2023}. TAID constructs a temporally adaptive intermediate distribution between student and teacher, gradually shifting the target as the student improves \cite{shingTAIDTemporallyAdaptive2025}. Selective token-weighted distillation masks or down-weights rejected tokens to avoid uniformly applying loss where the teacher signal is noisy or poorly matched to the student \cite{huang2025SelecTKDSelectiveTokenWeighted}. Speculative Knowledge Distillation lets the student propose tokens and the teacher replace low-quality choices, producing training data that is closer to the student's inference-time distribution \cite{xu2024SpeculativeKnowledgeDistillation}. Veto reformulates on-policy targets through a geometric bridge in logit space to suppress harmful gradients when teacher and student are too far apart \cite{jang2026StableOnPolicyDistillation}. In mathematical reasoning, the capacity gap is especially pronounced because small models may fail to benefit from long chain-of-thought traces produced by strong reasoners, and can perform better when trained on shorter or mixed-complexity reasoning traces aligned with their own capacity \cite{liSmallModelsStruggle2025}.

These methods show that the capacity gap is not only a question of teacher quality; it is also a question of whether the training setup exposes information that the student can actually use. Our approach is deliberately simpler. We do not add intermediate teachers, change the student architecture, construct adaptive target distributions, selectively mask loss terms, interleave teacher and student sampling, geometrically reformulate the target distribution, or shorten teacher reasoning traces. Instead, we intervene at the prompt level through context supplied to either the teacher or the student, then study whether the resulting behavior can be distilled into the trained student. That way we can keep the standard loss definitions and training loops. This connects the capacity-gap problem to context distillation, privileged context, and self-distillation.

\section{Context distillation and privileged context}

Context distillation refers to training a model so that behavior induced by an explicit context is internalized into the model parameters. The context can be an instruction, explanation, scratchpad, demonstration, optimized prompt, or other conditioning signal that improves behavior when it is present. The goal is to preserve the effect of that context after training, rather than requiring the full context to be supplied at inference time. Early alignment work introduced this idea for long behavioral prompts \cite{askell2021GeneralLanguageAssistant}, and subsequent work framed it more generally as distilling useful context tokens such as instructions, explanations, scratchpads, and examples into the model weights \cite{snell2022LearningDistillingContext}.

The most direct version is off-policy context distillation. In this setting, the context is associated with the teacher: the teacher receives both the original input and the additional context, produces rollouts under that conditioning, and the student is trained to imitate those rollouts. Knowledge-update distillation follows this pattern by prompting a model with new entity definitions and then training the unconditioned model to match the definition-conditioned behavior \cite{padmanabhan2023PropagatingKnowledgeUpdates}. Prompt Baking similarly turns a fixed behavioral prompt into a weight update by making the baked model match the distribution of the prompted model \cite{bhargava2024PromptBaking}. Generative Prompt Internalization argues that output imitation alone can be an indirect signal for learning the hidden prompt, and therefore also trains the model to generate prompt content and reasons for the behavioral change \cite{shin2025GenerativePromptInternalization}. These methods differ in how explicitly they expose or reconstruct the hidden context, but they share the same off-policy structure: the training trajectories are generated by a context-conditioned teacher, and the student learns the context only through imitation of those teacher-induced trajectories.

On-policy context distillation introduces the context in two different places. The first route gives the context to the teacher as privileged context. The student generates rollouts from its ordinary input distribution, while the teacher scores the student's prefixes with access to additional information unavailable to the student. The context is therefore transferred through teacher log-probabilities rather than through teacher-written trajectories. This is the setting studied by On-Policy Context Distillation, which trains the student on its own trajectories while minimizing a reverse-KL objective against a context-conditioned teacher \cite{yeOnPolicyContextDistillation2026}. Privileged-information distillation and related self-distillation methods use the same asymmetry in agentic or reasoning settings: $\pi$-Distill and On-Policy Self-Distillation condition the teacher on privileged action information \cite{penalozaPrivilegedInformationDistillation2026}, while Self-Distilled Reasoner uses a single model in two roles, with the teacher role conditioned on verified reasoning traces and the student role left unconditioned \cite{zhaoSelfDistilledReasonerOnPolicy2026}.

A broader self-distillation literature also converts richer feedback into dense token-level supervision by conditioning one role of the model on information unavailable to the other. Self-Distillation Fine-Tuning uses demonstration-conditioned in-context behavior as a teacher for continual learning \cite{shenfeldSelfDistillationEnablesContinual2026}. Self-Distillation Policy Optimization conditions on textual feedback such as runtime errors or judge evaluations and distills the feedback-informed predictions back into the policy \cite{hubotterReinforcementLearningSelfDistillation2026}. SD-Zero conditions a reviser on an initial response and binary reward, then distills the reviser into a generator \cite{heSelfDistillationZeroSelfRevision2026}. These methods show that privileged or feedback-conditioned predictions can provide dense supervision, but they usually study self-distillation or teacher--student pairs with relatively matched capacities. They also reveal risks: directly relying on privileged teachers conditioned on reference answers can lead to leakage and unstable training \cite{yangSelfDistilledRLVR2026,yangSelfDistilledRLVR2026a}, and self-distillation in reasoning can suppress epistemic verbalization and degrade out-of-distribution performance when the context removes useful uncertainty from the teacher behavior \cite{kimWhyDoesSelfDistillation2026}.

The second on-policy route gives the context directly to the student while it generates rollouts. In this setting, the rollout distribution itself is context-conditioned: the student sees the instruction or behavioral context during generation, and the teacher then scores prefixes produced by this context-aware student. This differs from privileged-context distillation because the context changes the states visited during training, not only the teacher probabilities assigned to those states. To the best of our knowledge, prior context-distillation work has not studied this student-visible on-policy variant as a way to train the model to internalize an instruction supplied during rollout generation in this capacity gap regime. Both on-policy routes aim to produce a trained student whose behavior incorporates the instruction encoded in the context: the privileged-teacher route transfers the instruction indirectly through teacher log-probabilities, while the student-visible route lets the instruction shape the student's own training trajectories.
